Until now $\X$ could be imprimitive.
Now we specialize to the primitive case to complete the proof of Theorem~\ref{main-thm}.
Let $\alpha \in \fS$ be minimal.
If $\X$ is primitive then $R_\alpha$ must be connected,
which implies that $\alpha^*$ covers $\{1, \dots, d\}$.
Hence $\alpha^*$ is an equipartition of~$\{1, \dots ,d\}$.
Let $w = \wt(\alpha)$.
If $w = 1$ then $\alpha^*$ must be the set of elements of weight $1$, so $R_\alpha$ is the Cameron graph.
Since the Weisfeiler--Leman stabilization of the Cameron graph is the Cameron scheme, we find $\X \le \cC(m, k, d)$.
If $w > 1$ then~$\alpha^* \subset \{0, k\}^d$.
If the elements of $\alpha^*$ have support size $e$ 
then $R_\alpha$ is the Hamming graph $H(M^e, d/e)$, so $\X \le \cH(d/e, M^e)$.
To finish we must show $\cT_{M^e}^{d/e} \le \X$.
For this it suffices to prove that for every $\beta \in \fS$, the elements of $\beta^*$ are sums of elements of $\alpha^*$ (and in particular~$\beta^* \subset \{0, k\}^d$).

We apply Proposition~\ref{key-prop}(2) to $\alpha$ and $\beta$.
Let $i \in \{1, \dots, d\}$.
Taking $a = e_i$, we find that $N^\beta_\alpha$ is at least the number of $b \in \beta^*$ such that $b_i > 0$.
On the other hand, taking~$a$ to be the unique element of $\alpha^*$ dominating $e_i$,
since $a \in \{0,k\}^d$ we find that~$N^\beta_\alpha$ is equal to the number of $b \in \beta^*$ such that $a \le b$.
Hence $b_i > 0$ implies $a \le b$.
This implies that $b$ is the sum of those $a \in \alpha^*$ such that $a \le b$, as required.
